\documentclass[10pt,a4paper]{amsart}
\usepackage[margin=19mm]{geometry}
\usepackage{amsmath,amssymb,mathtools}
\usepackage[scr=rsfs,cal=euler]{mathalfa}
\usepackage{xfrac}
\usepackage[unicode,hidelinks,bookmarks=false]{hyperref}
\newcommand{\ie}{i.e.\ }
\newcommand{\cf}{cf.\ }
\newcommand{\resp}{respectively\ }
\newcommand{\Subsection}{\subsection}
\providecommand{\nobreakdash}{\nobreak\hbox{-}}
\setlength{\emergencystretch}{2em}
\multlinegap0pt
\usepackage{bm}
\usepackage{bbm}
\usepackage{dsfont}
\usepackage{xspace}
\usepackage{cancel}
\usepackage{mathtools}
\usepackage{amsthm}
\makeatletter
\@ifundefined{lem}{\newtheorem{lem}{Lemma}}{}
\@ifundefined{prop}{\newtheorem{prop}{Proposition}}{}
\makeatother
\DeclareMathOperator{\ruc}{\xrightarrow[]{r}}
\renewcommand{\ge}{\geqslant}
\renewcommand{\le}{\leqslant}
\title[Free uniformly complete vector lattices]{Free uniformly complete vector lattices}
\author{Eduard Emelyanov, Svetlana Gorokhova}
\date{Source: arXiv:2109.03895v5 (CC BY 4.0)}
\begin{document}
\maketitle

\noindent\emph{Source and licence.} Free uniformly complete vector lattices. Version arXiv:2109.03895v5 (CC BY 4.0), distributed under the Creative Commons Attribution 4.0 International licence, as recorded in the arXiv licence metadata for this exact version and shown on the abstract page https://arxiv.org/abs/2109.03895v5. Full rights notice: licenses/arxiv-2109.03895-CC-BY-4.0.txt.\par\smallskip

\noindent\textit{Editorial scope.} Counted unit: the Prop whose source label is \texttt{L2} in the article (source0.tex lines 289--293, source label \texttt{L2}), delivered together with the article's own complete original proof of it (source0.tex lines 295--362, one \texttt{proof} environment of 68 lines). No mathematics is written, repaired or re-derived: statement and proof are the article's own text line for line. Required context retained uncounted: L1 (lines 260--263). Every definition, display and label of those lines is retained unchanged. Omitted from this package and not redistributed: 1--259, 264--288, 294--294, 363--657. Nothing inside a retained excerpt is omitted or altered: every retained body is delivered exactly as the article's lines are written, so no omission receipt of any kind accompanies this package. Citations: the retained excerpts carry no citation command at all, so no bibliography entry of the article is transcribed and no reference is redistributed. Reference adaptation: none was needed and none was made; every reference inside the delivered unit resolves inside it, so no unresolved reference or citation is produced. Printed slips kept literal: no printed word, symbol, hypothesis or number of the counted statement or of its proof was altered. Layout and class: the article's own class is \texttt{article} with packages including amsmath, amsthm, enumerate, amssymb, bbold, color, titlesec; the wrapper is the lane's pinned \texttt{amsart} editorial wrapper at 10pt on A4, which restates the theorem-like environments the retained text needs and adds the article's own macro definitions that the retained text needs (\texttt{\textbackslash ge}, \texttt{\textbackslash le}, \texttt{\textbackslash ruc}), with extra wrapper packages (bm, bbm, dsfont, xspace, cancel, mathtools). The delivered numbering is the unit's own. Assets: no retained span contains a figure environment, a graphics-inclusion command, a PDF-page inclusion, a table or a TikZ picture; no image file is copied into this package, the package declares no asset dependency and no image file is redistributed with it. Licence: version arXiv:2109.03895v5 of this article is distributed under the Creative Commons Attribution 4.0 International licence, as recorded in the arXiv licence metadata for this exact version and shown on the abstract page https://arxiv.org/abs/2109.03895v5. A full rights notice is delivered at licenses/arxiv-2109.03895-CC-BY-4.0.txt. Characters: every retained excerpt is pure ASCII, so the delivered engine, XeLaTeX with its default Latin Modern text font, reports no missing character.
\begin{lem}[]\label{L1}
The intersection $F$ of all \text{\rm UCVL}s of $U$ containing $E$ satisfies
$F=\bigcup_{\gamma\in{\text{\rm Ord;$\gamma<\omega_1$}}}F_{\gamma}$.
\end{lem}

\begin{prop}[]\label{L2}
Let $E$ be a \text{\rm VSL} of a \text{\rm UCVL} $U$. 
Then the intersection $F$ of all \text{\rm UCVL}s of $U$ containing $E$
is an $\text{\rm r}$-completion of $E$.
\end{prop}

\begin{proof}
By Lemma \ref{L1}, $F$ is a UCVL that coincides with 
$\bigcup_{\gamma\in{\text{\rm Ord;$\gamma<\omega_1$}}}F_{\gamma}$.
In order to show that $F$ is an $\text{\rm r}$-completion of $E$,
let $G$ be a UCVL and let $T: E\to G$ be a Riesz 
homomorphism. It is enough to show that $T$ has
a unique extension to a Riesz homomorphism $S: F\to G$. 

\smallskip
First we prove that such an extension $S$,
if exists, is unique. Indeed assume $S_1: F\to G$ is another such an extension.
Clearly, $S$ agrees with $S_1$ on $F_1=E$. 
Assume $S$ agrees with $S_1$ on $F_{\beta}$ for all $\beta<\gamma$.
Let $x\in F_{\gamma}$.
If $\gamma$ is a limit ordinal then
$x\in F_{\gamma}=\bigcup_{\beta<\gamma}F_{\beta}$.
Then $x\in F_{\beta}$ for some $\beta<\gamma$, and hence
$Sx=S_1x$ by the assumption.
If $\gamma=\beta+1$ then there exists a sequence 
$x_n$ in $F_\beta$ with $x_n \ruc x(x_1)$. In other words, for each $k\in\mathbb{N}$, 
there exists $n_k$ with $|x_n-x|\le\frac{1}{k}|x_1|$ for all $n\ge n_k$.
Then 
$$
   |Sx_n-Sx|+|S_1x_n-S_1x|\le\frac{1}{k}(S|x_1|+S_1|x_1|)\ \ \ \ \ (\forall n\ge n_k).
   \eqno(2)
$$
Because the sequence $x_n$ is contained in $F_\beta$ then $Sx_n=S_1x_n$ 
for all $n\in\mathbb{N}$. Now, it follows from (2) that
$$
   Sx_n\ruc Sx(S|x_1|+S_1|x_1|) \ \ \ \text{\rm and} \ \ \ 
   S_1x_n\ruc S_1x(S|x_1|+S_1|x_1|)
$$ 
in $G$, and since $G$ is Archimedean, then $Sx=S_1x$ as desired.

\smallskip
In order to show the existence of the required extension $S: F\to G$,
we use Lemma \ref{L1} and apply the transfinite induction by $\gamma$.
Suppose the extensions $S_\alpha: F_\alpha\to G$ exist already
for all $\alpha<\gamma$. By the arguments as above,
they are unique for all $\alpha<\gamma$,
and hence $S_\alpha$ agree with $S_\beta$ on $F_{\min{(\alpha,\beta)}}$.
Let $x\in F_\gamma$.

\smallskip
(A) If $\gamma$ is a limit ordinal then $x\in F_\alpha$ 
for some $\alpha<\gamma$. We set $S_\gamma x:=S_\alpha x$.

\smallskip
(B) If $\gamma=\beta+1$, there exists a sequence 
$x_n$ in $F_\beta$ with $x_n \ruc x(x_1)$. Then, for each $k\in\mathbb{N}$, 
there exists $n_k$ with $|x_n-x|\le\frac{1}{k}|x_1|$ for all $n\ge n_k$, and hence 
$$
   |x_n-x_m|\le|x_n-x|+|x_m-x|\le\frac{2}{k}|x_1| \ \ \ \ (\forall n,m\ge n_k).
   \eqno(3)
$$
It follows from (3) that $|S_\beta x_n-S_\beta x_m|\le\frac{2}{k}S_\beta|x_1|\in G$ 
for all $n,m\ge n_k$.
So, the sequence $S_\beta x_n$ is $\text{\rm r}$-Cauchy in $G$, and hence there exists 
unique $y\in G$ with $S_\beta x_n\ruc y$ in $G$. 
It is easy to see that $y$ does not depend on the choice of the sequence 
$x_n$ in $F_\beta$ such that $x_n \ruc x(x_1)$. We set $S_\gamma x:=y$.

\smallskip
Finally, let $x\in F$. Then $x\in F_\gamma$ for some $\gamma$. 
We set $Sx:=S_\gamma x$. Since linear and lattice operations in VLs
are $\text{\rm r}$-continuous, the constructed extension $S: F\to G$ 
is a Riesz homomorphism. 
\end{proof}

\end{document}
